%% Curriculum vitae of Abolfazl Ziaeemehr, built with moderncv.
%% Compile: tectonic Ziaeemehr.tex   (or xelatex / pdflatex)

\PassOptionsToPackage{colorlinks=true}{hyperref}
\documentclass[11pt,a4paper,sans]{moderncv}

\moderncvstyle{classic}
\moderncvcolor{blue}
% Match the website accent (teal).
\definecolor{color1}{HTML}{0B6B66}
\definecolor{color2}{HTML}{5A6870}

\usepackage[utf8]{inputenc}
\usepackage[scale=0.8, top=1.6cm, bottom=1.8cm]{geometry}
\setlength{\hintscolumnwidth}{2.9cm}

% personal data
\name{Abolfazl}{Ziaeemehr}
\title{Research engineer, scientific software and machine learning}
\address{Amiens}{France}
\phone[mobile]{+33~6~05~62~43~46}
\email{a.ziaee@u-picardie.fr}
\homepage{ziaeemehr.github.io}
\social[github]{Ziaeemehr}
\social[linkedin]{ziaeemehr}
\photo[72pt][0.4pt]{../images/myimage.png}

\newcommand{\me}{\textbf{A.~Ziaeemehr}}
\AtBeginDocument{\hypersetup{urlcolor=color1, linkcolor=color1, citecolor=color1}}

\begin{document}
\makecvtitle

\section{Profile}
\cvitem{}{I build fast, differentiable software for simulating dynamical systems and fitting them to data.
At Aix-Marseille University I led the development of VBI, a toolkit for simulation-based
inference on whole-brain models, and in 2026 I wrote and released JaxCont and lyapax (numerical
analysis in JAX) and LitForge (a local-first literature app). I care about tested code, documentation and
releases that other people can install.}

\section{Experience}
\cventry{Oct 2026--now}{Research engineer (Ingénieur de recherche)}{Université de Picardie Jules Verne}{Amiens, France}{}{}
\cventry{Oct 2021--Sep 2026}{Postdoctoral researcher}{Institut de Neurosciences des Systèmes, Inserm UMR 1106, Aix-Marseille University}{Marseille, France}{}{%
With Viktor Jirsa and Meysam Hashemi.
\begin{itemize}
\item Lead developer of VBI: C++, CuPy and Numba simulators, feature extraction and PyTorch neural density estimators for inferring brain-model parameters from EEG, MEG, SEEG and fMRI.
\item Applied simulation-based inference in studies of epilepsy, multiple sclerosis, Parkinson's disease, Alzheimer's disease and schizophrenia.
\item Contributor to the \textit{sbi} library; taught high-performance Python to the group.
\end{itemize}}
\cventry{Mar--Sep 2021}{Senior scientific developer}{Panoptopia}{}{}{Python packages for costing and risk management.}
\cventry{Sep 2020--Mar 2021}{Researcher}{Institute for Research in Fundamental Sciences (IPM)}{Tehran}{}{Spiking and rate models of the basal ganglia in Parkinson's disease.}
\cventry{Apr 2019--Feb 2020}{Research assistant}{IPM}{Tehran}{}{Chimera states and population models on the human connectome.}
\cventry{Apr--Sep 2018}{Visiting researcher}{University of Granada, Computational Physics Group}{Spain}{}{Phase transitions and noise-induced volatility in networks of integrate-and-fire neurons.}

\section{Software}
\cvitem{VBI}{\href{https://github.com/ins-amu/vbi}{Virtual Brain Inference}, lead developer. Simulation-based inference for whole-brain models; published in eLife (2025) and JOSS (2026). PyPI, EBRAINS.}
\cvitem{JaxCont}{\href{https://github.com/Ziaeemehr/JaxCont}{Differentiable continuation and bifurcation analysis in JAX}: equilibrium and periodic-orbit branches, fold, Hopf, branch-point and period-doubling detection, two-parameter curves; compiles with \texttt{jit}, batches with \texttt{vmap}, differentiates with \texttt{grad}. Validated against \mbox{MatCont~7.6}. PyPI, Read the Docs, Zenodo.}
\cvitem{lyapax}{\href{https://github.com/Ziaeemehr/lyapax}{Lyapunov exponent spectra of ODEs, maps and delay equations in JAX} via forward-mode tangent propagation; runs on GPU; benchmarked against ChaosTools.jl. PyPI, Read the Docs.}
\cvitem{LitForge}{\href{https://github.com/Ziaeemehr/LitForge}{Local-first desktop app} for literature discovery and screening, question answering with verified quotes, and LaTeX writing with local or cloud language models.}
\cvitem{MD Viewer}{\href{https://github.com/Ziaeemehr/MDviewer}{Bidirectional Persian and English Markdown viewer} with KaTeX, released for web, desktop (Tauri) and Android.}
\cvitem{Other}{Contributions to sbi, NESTML and Brian2; JITCSIM, vbjax, mndynamics; reproducible code for each of my first-author papers. Full list at \href{https://ziaeemehr.github.io/projects.html}{ziaeemehr.github.io}.}

\section{Skills}
\cvitem{Languages}{Python, C++, JavaScript, Julia, Bash}
\cvitem{Numerics}{JAX, NumPy, SciPy, Numba, CuPy and CUDA; ODE, SDE and DDE solvers, continuation and bifurcation methods, Lyapunov analysis}
\cvitem{Machine learning}{PyTorch, simulation-based inference (sbi), normalising flows and neural density estimators, CNNs for EEG}
\cvitem{Data}{EEG, MEG, SEEG and fMRI; MNE-Python, The Virtual Brain, NEST, Brian2}
\cvitem{Engineering}{Git, GitHub Actions, Linux, PyPI packaging and releases, pytest, Sphinx and Read the Docs, SWIG bindings, Tauri, HPC clusters}

\section{Education}
\cventry{2014--2020}{PhD, computational neuroscience}{Institute for Advanced Studies in Basic Sciences (IASBS)}{Zanjan, Iran}{}{Thesis: \textit{Synchronization dynamics on undirected and directed hierarchical networks}. Supervisors: Mina Zarei and Alireza Valizadeh.}
\cventry{2008--2011}{MSc, solid-state physics}{Imam Khomeini International University}{Qazvin, Iran}{}{Thesis: \textit{Calculation of optical properties of metallic nanoparticles}.}
\cventry{2004--2008}{BSc, physics}{University of Qom}{Iran}{}{}

\section{Publications}
\cvitem{2026}{G. Preti, H. E. Wang, \me, et al. \textbf{Virtual brain twins guide personalized treatment decision in schizophrenia}. \href{https://doi.org/10.64898/2026.05.06.26352533}{\textit{medRxiv}}, 2026.}
\cvitem{2026}{A. Esmaeili, M. Woodman, N. Baldy, \me, et al. \textbf{Cohort-amortized personalization: navigating the privacy-utility frontier for virtual brain twins}. \href{https://doi.org/10.48550/arXiv.2606.30329}{\textit{arXiv}}, 2026.}
\cvitem{2026}{\me, L. Domide, M. Woodman, S. Petkoski, V. Jirsa, M. Hashemi. \textbf{Virtual Brain Inference (VBI) toolkit on EBRAINS}. \href{https://doi.org/10.21105/joss.10612}{\textit{Journal of Open Source Software}} 11(124), 10612, 2026.}
\cvitem{2026}{M. Hashemi, D. Depannemaecker, M. Saggio, \ldots, \me, et al. \textbf{Principles and operation of virtual brain twins}. \href{https://doi.org/10.1109/RBME.2025.3562951}{\textit{IEEE Reviews in Biomedical Engineering}} 19, 111--139, 2026.}
\cvitem{2025}{\me, M. Woodman, L. Domide, S. Petkoski, V. Jirsa, M. Hashemi. \textbf{Virtual Brain Inference (VBI), a flexible and integrative toolkit for efficient probabilistic inference on whole-brain models}. \href{https://doi.org/10.7554/eLife.106194}{\textit{eLife}} 14, RP106194, 2025.}
\cvitem{2025}{A. Esmaeili, M. Demolliens, M. Viersen, \me, et al. \textbf{Probabilistic inference of social presence across brain scales reveals enhanced synaptic efficacy}. \href{https://doi.org/10.1038/s42003-025-08971-3}{\textit{Communications Biology}} 8, 1608, 2025.}
\cvitem{2025}{J. Boelts, M. Deistler, M. Gloeckler, \ldots, \me, \ldots, J. H. Macke. \textbf{sbi reloaded: a toolkit for simulation-based inference workflows}. \href{https://doi.org/10.21105/joss.07754}{\textit{Journal of Open Source Software}} 10(108), 7754, 2025.}
\cvitem{2025}{G. Rabuffo, H. Lokossou, Z. Li, \me, et al. \textbf{Mapping global brain reconfigurations following local targeted manipulations}. \href{https://doi.org/10.1073/pnas.2405706122}{\textit{PNAS}} 122(16), e2405706122, 2025.}
\cvitem{2025}{C. Mazzara, \me, E. Troisi Lopez, et al. \textbf{Mapping brain lesions to conduction delays: the next step for personalized brain models in multiple sclerosis}. \href{https://doi.org/10.1002/hbm.70219}{\textit{Human Brain Mapping}} 46(7), e70219, 2025.}
\cvitem{2025}{M. Angiolelli, D. Depannemaecker, H. Agouram, \ldots, \me, et al. \textbf{The virtual parkinsonian patient}. \href{https://doi.org/10.1038/s41540-025-00516-y}{\textit{npj Systems Biology and Applications}} 11, 40, 2025.}
\cvitem{2024}{M. Hashemi, \me, M. Woodman, J. Fousek, S. Petkoski, V. Jirsa. \textbf{Simulation-based inference on virtual brain models of disorders}. \href{https://doi.org/10.1088/2632-2153/ad6230}{\textit{Machine Learning: Science and Technology}} 5(3), 035019, 2024.}
\cvitem{2024}{P. Sorrentino, A. Pathak, \me, et al. \textbf{The virtual multiple sclerosis patient}. \href{https://doi.org/10.1016/j.isci.2024.110101}{\textit{iScience}} 27(7), 110101, 2024.}
\cvitem{2023}{B. H. Yalçınkaya, \me, J. Fousek, et al. \textbf{Personalized virtual brains of Alzheimer's disease link dynamical biomarkers of fMRI with increased local excitability}. \href{https://doi.org/10.1101/2023.01.11.23284438}{\textit{medRxiv}}, 2023.}
\cvitem{2021}{\me, A. Valizadeh. \textbf{Frequency-resolved functional connectivity: role of delay and the strength of connections}. \href{https://doi.org/10.3389/fncir.2021.608655}{\textit{Frontiers in Neural Circuits}} 15, 608655, 2021.}
\cvitem{2020}{\me, M. Zarei, A. Valizadeh, C. R. Mirasso. \textbf{Frequency-dependent organization of the brain's functional network through delayed-interactions}. \href{https://doi.org/10.1016/j.neunet.2020.08.003}{\textit{Neural Networks}} 132, 155--165, 2020.}
\cvitem{2020}{\me, M. Zarei, A. Sheshbolouki. \textbf{Emergence of global synchronization in directed excitatory networks of type I neurons}. \href{https://doi.org/10.1038/s41598-020-60205-0}{\textit{Scientific Reports}} 10, 3306, 2020.}

\section{Teaching}
\cvitem{2026}{\href{https://github.com/Ziaeemehr/dl_eeg}{Deep learning by doing}: self-study course from tensors to CNNs for EEG classification.}
\cvitem{2024}{\href{https://github.com/Ziaeemehr/workshop_hpcpy}{High-performance computing with Python}, workshop for the Theoretical Neuroscience Group, INS.}
\cvitem{Jul 2020}{Teaching assistant, Neuromatch Academy (three-week summer school).}
\cvitem{2016--2017}{Workshop lecturer at IASBS: Python scripting, Julia, C++, Brian2 and NEST.}
\cvitem{2015--2016}{Teaching assistant: classical electrodynamics I and II, computational physics.}

\section{Awards}
\cvitem{2018}{Ministry of Science of Iran scholarship for research at the University of Granada, Spain.}
\cvitem{2014}{Ranked 26th of about 5,000 candidates in the national PhD entrance exam.}

\section{Languages}
\cvitemwithcomment{Persian}{Native}{}
\cvitemwithcomment{English}{Fluent}{}
\cvitemwithcomment{French}{A2}{}

\section{References}
\cvitem{}{Viktor Jirsa, Director of Research, INS, Aix-Marseille University\newline \href{mailto:viktor.jirsa@univ-amu.fr}{viktor.jirsa@univ-amu.fr}}
\cvitem{}{Meysam Hashemi, Senior Researcher, INS, Aix-Marseille University\newline \href{mailto:meysam.hashemi@univ-amu.fr}{meysam.hashemi@univ-amu.fr}}
\cvitem{}{Alireza Valizadeh, Professor of Physics, IASBS\newline \href{mailto:valizade@iasbs.ac.ir}{valizade@iasbs.ac.ir}}
\cvitem{}{Mina Zarei, Associate Professor of Physics, IASBS\newline \href{mailto:mina.zarei@iasbs.ac.ir}{mina.zarei@iasbs.ac.ir}}

\end{document}
